\chapter{More About Free Resolutions}
The goal of this appendix is to prove
\Cref{thm:tor_properties} (a),
that $\Tor_i^A(M,N)$ 
does not depend on the choice of free resolution of $N$.
To do this, we need to introduce a little more 
homological algebra.
\section{Chain Homotopy}
First, we define morphism between two chain complexes.
\begin{definition}[Morphism between chain complexes]
    Let $C_*$ and $D_*$ be two chain complexes of $A$-modules.
    A \vocab{morphism} between two chain complexes
    $f : C_*\to D_*$ is the sequence of morphisms
    $\dots, f_0, f_1,\dots$ such that the following diagram commutes
    \begin{center}
    \begin{tikzcd}
        \dots \arrow[r] 
        & C_1 \arrow[d, "f_{-1}"] \arrow[r, "\partial_1"]
        & C_0 \arrow[d, "f_0"] \arrow[r, "\partial_0"]
        & C_{-1} \arrow[d, "f_1"] \arrow[r]
        & \dots  \\
        \dots \arrow[r] 
        & D_1 \arrow[r, "\partial_1"]
        & D_0 \arrow[r, "\partial_0"]
        & D_{-1} \arrow[r]
        & \dots   
    \end{tikzcd}
    \end{center}
    This promotes the set of chain complexes
    of $A$-modules to be a category $\operatorname{ch}(\Mod A)$.
    \notation{$\operatorname{ch}(\Mod A)$}{category
    of chain complexes of $A$-modules}
\end{definition}
For any morphism $f : C_* \to D_*$,
you can check that $f_n$ induces a map
$H_n(f) : H_n(C_*) \to H_n(D_*)$.
(This makes $H_n$ a functor 
$H_n : \operatorname{ch}(\Mod A) \to \Mod A$.)
\notation{$H_n(f)$}{induced map of homology groups}
For any two maps $f,g$, we are interested
in the condition at which $H_n(f) = H_n(g)$.
One sufficient condition is chain homotopy.
\begin{definition}[Chain homotopy]
    Let $C_*$ and $D_*$ be two chain complexes of $A$-modules.
    A \vocab{chain homotopy} between two morphisms
    $f,g : C_*\to D_*$ is a sequence of morphism 
    $h_n : C_n \to D_{n+1}$ such that
    $$\partial_{n+1}h_n - h_{n-1}\partial_n = g_n-f_n.$$
    If there exists such $h$, then we say that
    $f$ and $g$ are \vocab{chain homotopic}.
\end{definition}
\begin{proposition}
    \label{prop:chain_homotopy}
    If $h$ is a chain homotopy between
    $f,g : C_*\to D_*$, then $H_n(f) = H_n(g)$
    (as homomorphisms from $H_n(C_*)\to H_n(D_*)$) for all $n$.
\end{proposition}
\begin{proof}
    Suppose that $z\in \Ker(\partial_n)$,
    i.e., $\partial_n(z)=0$.
    We need to show that $f(z)-g(z) \in \Im(\partial_{n+1})$.
    To do this, we note that
    \[g_n(z)-f_n(z) = \partial_{n+1}(h_n(z)) - h_{n-1}(\partial_n(z))
    =\partial_{n+1}(h_n(z)) \in \Im\partial_{n+1}
    \qedhere\]
\end{proof}

\section{Fundamental Theorem of Homological Algebra}
\begin{theorem}[Fundamental theorem of homological algebra]
\label{thm:homological_alg}
Let $N$ and $N'$ be two $A$-modules.
Let $N_*$ and $N'_*$ be free resolutions of $N$ and $N'$.
Then, any morphism $\phi : N\to N'$ can be lifted to
morphism between free resolutions $\widetilde\phi : N_*\to N'_*$.
Furthermore, any two lifts are chain homotopic.
\end{theorem}
\begin{proof}
    We use the following properties labeled ($\spadesuit$) 
    of free modules:
    if $F$ is a free module,
    then for any surjective map $f : P\to Q$,
    it is possible to lift any map $\pi : F\to Q$ to 
    $\widetilde\pi : F\to P$
    \begin{center}
    \begin{tikzcd}
        & F \arrow[ld, dashed, "\exists\widetilde\pi", swap] 
        \arrow[d, "\psi"]\\
        P \arrow[r, two heads, "f", swap] & Q
    \end{tikzcd}.
    \end{center}
    (Aside: Modules with this property are called 
    \vocab{projective modules}, and this is the only condition
    of $N_i$ that we need to make this theorem works,
    so instead of free resolutions, you can use a more general projective resolutions.)
    To verify this property, simply send
    a generator $e_i\in F$ to any element in 
    $f^{-1}(\{\psi(e_i)\})$.

    Now, we prove the theorem. 
    
    \textbf{Existence of a lift.}
    First, we show existence
    by construction the map $\widetilde\phi_i : N_i\to N'_i$ by induction.
    In the base case, consider the following diagram
    \begin{center}
    \begin{tikzcd}
        N_0  \arrow[d, dashed] \arrow[r, two heads]
        \arrow[rd]
        & N  \arrow[d, "\phi"] \\
        N'_0 \arrow[r, two heads] & N' 
    \end{tikzcd}
    \end{center}
    Thus, we use $(\spadesuit)$ to find the map
    $\widetilde\phi_0$ such that the diagram above commute.

    Now, for the inductive step, suppose we have
    $\widetilde\phi_{i-1} : N_{i-1}\to N'_{i-1}$.
    Let $K_{i-1} = \Ker(N_{i-1}\to N_{i-2})$
    and $K'_{i-1} = \Ker(N'_{i-1}\to N'_{i-2})$
    Then, $\widetilde\phi_{i-1}$ 
    induces the map $\widetilde\phi_{i-1} : 
    K_{i-1} \to K'_{i-1}$. Then, we have the diagram
    \begin{center}
    \begin{tikzcd}
        N_i \arrow[d, dashed] \arrow[r, two heads]
        \arrow[rd]
        & K_{i-1}  \arrow[d, "\widetilde\phi_{i-1}"] \\
        N'_i \arrow[r, two heads] & K'_{i-1}
    \end{tikzcd},
    \end{center}
    and again, we can use $(\spadesuit)$
    to find $\widetilde\phi_i : N_i\to N'_i$
    so that the above diagram commutes.

    \textbf{Any two lifts are chain homotopic.}
    We have to construct the chain homotopy.
    Consider two lifts $\psi, \pi : N_*\to N_*'$.
    Let $\ell = \psi-\pi$.
    We need $h$ such that $\partial h - h\partial = \ell$.
    We, again, do this by induction. 
    Ths base case is to find $h_0$ so that the 
    following diagram commutes:
    \begin{center}
    \begin{tikzcd}
        &
        & N_0  \arrow[d, "\ell_0"]
        \arrow[lld, dashed, "h_0", swap]
        \arrow[ld, dashed] \arrow[r, two heads]
        & N \arrow[d, "0"] \\
        N'_1 \arrow[r, two heads] 
        & K'_1 \arrow[r, hook] 
        & N'_0  \arrow[r, two heads]
        & N'
    \end{tikzcd}
    \end{center}
    This is done in two steps.
    First, note that the composition 
    $N\to N_0'\to N'$ is zero map,
    so that composition factors through $K_1'$.
    Therefore, we get $N_0\to K_1'$.
    Then, use $(\spadesuit)$ again to find $N_0\to N_1'$.
    Since the term $h\partial$ vanishes,
    we get that the chain homotopy condition
    is verified at $0$.

    Now, the inductive step is similar. We have the diagram
    \begin{center}
    \begin{tikzcd}
        &
        & N_i 
        \arrow[lld, dashed, "h_i", swap]
        \arrow[ld, dashed] 
        \arrow[d, "\ell_i + h_{i-1}\partial_i"] \\
        N'_{i+1} \arrow[r, two heads] 
        & K'_{i+1} \arrow[r, hook] 
        & N'_i  
    \end{tikzcd}
    \end{center}
    We have to construct $h_i$ such that
    $\partial_{i+1}h_i - h_{i-1}\partial_i = \ell_i$.
    To do this, we first note that
    \begin{align*}
    \partial_i\Big(\ell_i + h_{i-1}\partial_i\Big)
    &= \partial_i\ell_i + \partial_i h_{i-1}\partial_i \\
    &= \ell_{i-1}\partial_i + \partial_ih_{i-1}\partial_i 
    \tag{$\ell$ commutes with boundary maps}\\
    &= (\ell_{i-1} + \partial_{i}h_{i-1})\partial_i \\
    &= h_{i-2}\partial_{i-1}\partial_i 
    \tag{chain homotopy definition}\\
    &= 0,
    \end{align*}
    so $\ell_i + h_{i-1}\partial_i$ factors through
    $K'_{i+1}$, getting the map 
    $N_i\to K_{i+1}'$. Finally, use $(\spadesuit)$
    to lift it to $h_i$, giving the desired
    chain homotopy condition.
\end{proof}
\section{Proof that Tor does not Depend on Free Resolution}
We now prove \Cref{thm:tor_properties} (a).
\begin{theorem}[Tor does not depend on free resolution]
    \label{thm:tor_free_res}
    The functor $\Tor_n^A(M,N)$ 
    (defined in \Cref{def:tor}) does not depend on the 
    choice of $N_*$.
\end{theorem}
\begin{proof}
    Consider two free resolutions $N_*$ and $N'_*$
    for $N$. Thus, by the fundamental theorem of homological algebra (\Cref{thm:homological_alg}),
    the identity map $\id : N\to N$ lifts to 
    $\phi : N_* \to N'_*$ and $\psi : N'_*\to N_*$.
    Then, $\phi\circ\psi : N_*\to N_*$
    is a lift of $\id$ as well. However, the identity map
    is one such liftlift of $\id$, so $\phi\circ\psi$
    is chain homotopic to identity.
    
    Tensoring with $M$, we get two maps
    $\phi_* : M\otimes N_* \to M\otimes N_*'$
    and $\psi_* : M\otimes N'_* \to M\otimes N_*$.
    Since everything is an additive functor,
    the chain homotopic condition is preserved through
    tensoring with $M$. 
    Thus, $\phi_*\circ\psi_*$ is chain homotopic to identity.
    Similarly, $\psi_* \circ\phi_*$ is chain homotopic to identity.
    Thus, by using \Cref{prop:chain_homotopy},
    $H_n(\phi_*)$ and $H_n(\psi_*)$
    provide isomorphism between
    $H_n(M\otimes N_*) \simeq H_n(M\otimes N'_*)$.
    This implies that $\Tor_n^A(M,N)$ gives the same result
    for two resolutions $N_*$ and $N'_*$.
\end{proof}
